\section{监控与评估指引}
\subsection{关键指标}
自我提升训练必须持续监测以下指标，以便在 win rate plateau、hallucination 走高或 inference compute 激增时迅速干预。

\begin{table}[H]
\centering
\begin{tabular}{p{5cm}p{10cm}}
\toprule
指标 & 说明 \\
\midrule
Win rate delta & 以 Alpaca Eval 2 / DeepSeek benchmarks 为基准，观察 verified response 与 draft 的 win rate 变化。 \\
Reward variance & Self-rewarding 中 reward score 的标准差若急剧上升，说明模型可能在 overthinking 或 hallucinating。 \\
Hallucination frequency & 用 SelfT evaluator 标注经过验证的 hallucination 样本率，在超过门槛时降低 cross-check depth。 \\
Inference compute & 评估 verified response 产生的额外 compute（token count × number of cross-checks），防止 evaluation 成本超过预算。 \\
\bottomrule
\end{tabular}
\caption{监控 self-rewarding 训练的核心指标。}
\end{table}

\begin{importantbox}{指标联动建议}
在 win rate 停滞时，优先查看 reward variance 与 hallucination frequency：若 hallucination 占比提高，先回退到 Self-Instruct seeds；若 reward variance 走高，则依赖 meta judge calibration 及时收敛。
\end{importantbox}

\subsection{自动化告警与反馈}
对跛脚指标设置阈值后，可以用 automated scripts 生成 alert：例如 reward variance 超过 0.3 就暂停迭代、hallucination frequency 连续三次 > 5\% 就触发 human-in-the-loop 抽查。可以结合 Self-Feeding 流程，把 alert 推送到 pipeline dashboard，使操作团队能在 `self-improvement` 中进行快速 close loop。

\begin{knowledgebox}{告警策略模板}
1. 设定频率：每轮 iteration 结束后立即计算指标。\newline
2. 分级响应：warning 触发 retry prompt；alert 触发 additional verified response。\newline
3. 归档日志：hallucination 样本与 reward 轨迹一同写入 `audit.csv` 供 future meta judge retraining。
\end{knowledgebox}

\subsection{本章小结}
有了这套指标与告警，团队可以透明追踪 self-rewarding 在实际训练过程中的状态，一旦出现 drift 即可采用 meta judge 重新校准或 self-instruct prompt replay。

